\subsection{Use of AI/LLM tools}
\label{app:ai-tools}

% TODO(author): this appendix was drafted by Claude from what it did in the
% Cowork session of 23 September 2026 and from the earlier version of this
% appendix. Check every level against what YOU did, and complete the rows marked
% TODO; the LLM-assisted docstrings in the code must tell the same story.

Large language models were used throughout this project, following the course
guidelines\footnote{\url{https://github.com/EducationalMaterialUiO/MachineLearningUiO/blob/main/LLM_Usage_Declaration_Guidelines.md}},
which also define the contribution levels below. All code was tested, all
derivations were checked, and every claim in the report was read, understood
and verified by the author.

\subsubsection{Tools used}
\begin{itemize}
    \item Anthropic Claude via \url{https://claude.ai}, Claude Code, and the Claude
    desktop app (cowork), September--October 2026
    (Claude Sonnet 5.0 and Opus 5.5, 5, and 4.8 models).
    \item OpenAI ChatGPT via \url{https://chatgpt.com}, September 2026.
\end{itemize}


\subsubsection{Text writing and editing}
\Cref{tab:llm-writing-contribution} gives the contribution level for each part
of the report.

\begin{table}[htbp]
    \centering
    \caption{LLM contribution to text writing and editing (levels 0--3).}
    \label{tab:llm-writing-contribution}
    {\footnotesize
    \begin{tabularx}{\columnwidth}{@{}lcX@{}}
        \toprule
        \textbf{Section} & \textbf{Level} & \textbf{Notes} \\
        \midrule
        Abstract & 3 & Drafted by Claude from the finished results; revised by the author. \\
        Introduction & 2 & Author's text; last three paragraphs revised by Claude. \\
        Methods & 2 & Author's text; corrected by Claude. \\
        Implementation & 2 & Author's text; revised by Claude. \\
        Results and discussion & 2 & Author's text; revised and improved by Claude. \\
        Conclusion & 2 & Drafted by the author; revised and imporved by Claude. \\
        Appendices & 3 & Drafted by Claude; revised by the author. \\
        \bottomrule
    \end{tabularx}}
\end{table}

\subsubsection{Code generation and assistance}
\Cref{tab:code-generation} gives the level for each source file (levels 0--4).
Functions with level 2 or above carry an \texttt{LLM-assisted} tag in their
docstring. Claude Code was also used to write docstrings and unit tests across
the code base.

\begin{table*}[htbp]
    \centering
    \caption{LLM contribution to the code (levels 0--4).}
    \label{tab:code-generation}
    {\footnotesize
    \begin{tabularx}{\textwidth}{@{}l c X@{}}
        \toprule
        \textbf{File} & \textbf{Level} & \textbf{Description} \\
        \midrule
        \texttt{design\_matrix.py}, \texttt{runge.py}, \texttt{metrics.py} & 0 & \\
        \texttt{ordinary\_least\_squares.py} & 1 & Claude suggested the singular-value cutoff of the pseudoinverse. \\
        \texttt{ridge.py}, \texttt{cost.py}, \texttt{shrinkage.py} & 0 & \\
        \texttt{base.py} & 2 & scikit-learn estimator interface. \\
        \texttt{autodiff.py} & 2 & JIT compilation. \\
        \texttt{optimizers.py} & 3 & Provided the framework of the \texttt{Optimizer} class and shape validation. \\
        \texttt{schedules.py} & 4 & Learning-rate schedules; reviewed and tested by us. \\
        \texttt{gradient\_descent.py} & 2 & Learning-rate scheduling. \\
        \texttt{degree\_sweep.py} & 3 & Refactored from our notebooks and generalized. \\
        \texttt{lasso.py} & 0 & \\
        \texttt{bootstrap.py} & 2 & Bias against the known $f$ added by Claude. \\
        \texttt{cross\_validation.py} & 2 & Per-fold output added by Claude. \\
        \texttt{benchmark.py} & 4 & Iterations-to-tolerance benchmark, written by Claude. \\
        \texttt{scripts/} & 3 & Experiments, figures and configuration, written in part by Claude,
                                inspired by the notebooks. \\
        \texttt{notebooks/} & 2 & Assistanse with plotting, providing feedback, and debugging. \\
        \texttt{tests/} & 3 & Tests for the new code written by Claude; others with Claude Code. \\
        \bottomrule
    \end{tabularx}}
\end{table*}

\subsubsection{Declared and undeclared use}
An example of a use that must be declared is the benchmark function
\texttt{iterations\_to\_tolerance} in \texttt{benchmark.py}: Claude wrote the
function, which appears in our submission, so it is declared here and in its
docstring. An example of a use that need not be declared is using an LLM as a tutor, for
instance to explain the difference between a gradient and a subgradient, and
then writing about it in our own words. % TODO(author): replace with an example
% from your own work (the text of this report was drafted with Claude, so it
% cannot serve as the undeclared example).
